\documentclass[11pt,a4paper]{article}
\usepackage[margin=25mm]{geometry}
\usepackage{fontspec}
\setmainfont{TeX Gyre Pagella}
\usepackage{amsmath,amssymb,amsthm,mathtools}
\usepackage{enumitem,microtype,xcolor}
\usepackage[colorlinks=true,linkcolor=blue!45!black,citecolor=blue!45!black,urlcolor=blue!45!black]{hyperref}
\usepackage{fancyhdr}
\pagestyle{fancy}
\fancyhf{}
\fancyhead[L]{\small Boundary and iterated flows}
\fancyhead[R]{\small Companion to the nonboundary note}
\fancyfoot[C]{\thepage}
\setlength{\headheight}{14pt}
\setlength{\emergencystretch}{3em}
\setlength{\parskip}{3pt}
\setlist{itemsep=3pt,topsep=5pt}
\allowdisplaybreaks[2]
\newcommand{\C}{\mathbb C}
\newcommand{\Z}{\mathbb Z}
\newcommand{\N}{\mathbb N}
\newcommand{\slc}{\mathfrak{sl}_2}
\newcommand{\slh}{\widehat{\mathfrak{sl}}_2}
\newcommand{\slt}{\widetilde{\mathfrak{sl}}_2}
\newcommand{\ad}{\operatorname{ad}}
\newcommand{\Ann}{\operatorname{Ann}}
\newcommand{\Hom}{\operatorname{Hom}}
\newcommand{\Span}{\operatorname{span}}
\newcommand{\I}{\mathcal I}
\newcommand{\X}{\mathcal X}
\newcommand{\D}{\mathcal D}
\newcommand{\T}{\mathcal T}
\newtheoremstyle{plainupright}{6pt}{5pt}{\normalfont}{}{\bfseries}{.}{.5em}{}
\theoremstyle{plainupright}
\newtheorem{proposition}{Proposition}[section]
\newtheorem{lemma}[proposition]{Lemma}
\newtheorem{theorem}[proposition]{Theorem}
\newtheorem{corollary}[proposition]{Corollary}
\newtheorem{remark}[proposition]{Remark}
\hypersetup{pdftitle={Boundary Localization, Simplicity, and Iterated Flows},
 pdfsubject={Results of the living note v8 not contained in the nonboundary note}}
\title{\vspace{-14mm}Boundary Localization, Simplicity, and Iterated Flows\\
for Smooth Affine \(\slh\)-Modules\\[3pt]
\large Companion to the nonboundary note: results of the living note v8 (2026-09-02)\\
not contained in \emph{nonboundary\_localization\_results} (2026-09-11)}
\author{}
\date{September 2026}
\begin{document}
\maketitle
\vspace{-9mm}
\begin{abstract}
The nonboundary note \cite{Nonboundary} studies the interior quotients
\(T_{f_c}M_{a,b}\), \(c<b\), in depth, and uses the boundary results only as a
black box. The present note collects the boundary-side results of the living
note v8 \cite{LivingV8} that are \emph{not} contained in \cite{Nonboundary}:
the detailed boundary isomorphism with its right-normal PBW base comparison,
the simplicity flow through the FGXZ criterion, the Bernstein--Sato-type
\(b\)-function, the full trichotomy for arbitrary modes, the version with the
degree derivation, the iterated boundary chains, the non-critical
\(\Omega\)-quotient flow, and the critical-level flow with preserved
\(\theta\)-parameters. All statements are proved directly; this note is not a
claim of independent peer review or formal verification.
\end{abstract}

\section{Setup and background}\label{sec:setup}

All vector spaces and tensor products are over \(\C\). Let
\[
\slh=(\slc\otimes\C[t,t^{-1}])\oplus\C K,\qquad \slt=\slh\oplus\C d,
\]
with \(x_r=x\otimes t^r\) and the normalization
\begin{align*}
[e_r,f_s]&=h_{r+s}+r\delta_{r+s,0}K,&
[h_r,h_s]&=2r\delta_{r+s,0}K,\\
[h_r,e_s]&=2e_{r+s},&
[h_r,f_s]&=-2f_{r+s},&
[d,x_r]&=rx_r.
\end{align*}
The \(e\)'s commute with one another, the \(f\)'s commute with one another,
and \(K\) is central. Fix \(a,b\in\Z\) with \(a+b\ge0\), set \(L=a+b+1\ge1\),
and put
\[
S_{a,b}=\Span\{K,\ h_i\ (i\ge0),\ e_i\ (i>a),\ f_j\ (j>b)\}.
\]
For a character \(\varphi:S_{a,b}\to\C\) write
\[
M^c_{a,b}(\varphi)=\Ind_{S_{a,b}}^{\slh}\C_\varphi,
\qquad
M^f_{a,b}(\varphi)=\Ind_{S_{a,b}}^{\slt}\C_\varphi,
\qquad u=1\otimes1_\varphi,
\]
and
\[
\kappa=\varphi(K),\qquad \lambda_i=\varphi(h_i),
\qquad \lambda_\partial:=\lambda_L=\varphi(h_L).
\]
The character condition forces
\[
e_iu=0\ (i>a),\qquad f_ju=0\ (j>b),\qquad
h_iu=\lambda_iu,\qquad \lambda_i=0\ (i>L).
\]
For an injective real root vector \(F\) on a module \(X\), set
\[
\D_FX=\D_FU(\slh)\otimes_{U(\slh)}X,\qquad
\T_FX=\D_FX/X,
\]
where \(\D_FU(\slh)=\{F^m:m\ge0\}^{-1}U(\slh)\). We shall use the FGXZ
irreducibility criterion \cite{FGXZ}:

\begin{theorem}[FGXZ criterion]\label{thm:fgxz}
The module \(M^c_{N_1,N_2}(\chi)\) is simple if and only if
\[
\chi\bigl(h(N_1+N_2+1)\bigr)\ne0\qquad\text{and}\qquad \chi(K)\ne-2.
\]
\end{theorem}

\section{The right-normal localized PBW basis}\label{sec:pbw}

The following normal form is the engine of the boundary isomorphism. It is the
right-normal companion of the left-normal PBW model used in \cite{Nonboundary}.

\begin{lemma}[Right-normal localized PBW]\label{lem:right-pbw}
Let \(F=f_r\) with \(r\le b\). Choose a PBW order for \(M^c_{a,b}(\varphi)\) in
which \(f_r\) is the \emph{last} free variable, and let \(\X_r\) be the
ordered monomials in the remaining free families
\(f_j\ (j\le b,\ j\ne r)\), \(e_i\ (i\le a)\), \(h_j\ (j<0)\). Then
\[
\{XF^qu:X\in\X_r,\ q\in\Z_{\ge0}\}
\]
is a basis of \(M^c_{a,b}(\varphi)\),
\[
\{XF^qu:X\in\X_r,\ q\in\Z\}
\]
is a basis of \(\D_FM^c_{a,b}(\varphi)\), and
\[
\{XF^{-m}u+M^c_{a,b}(\varphi):X\in\X_r,\ m\ge1\}
\]
is a basis of \(\T_FM^c_{a,b}(\varphi)\). Correspondingly,
\[
\{d^kXF^qu:k\ge0,\ X\in\X_r,\ q\in\Z\}
\]
is a basis of \(\D_FM^f_{a,b}(\varphi)\), and in the quotient only \(q<0\)
remains.
\end{lemma}

\begin{proof}
The operator \(\ad F\) is locally nilpotent on \(U(\slh)\) and \(U(\slt)\), and
for the derivation \([F,d]=-rF\), so \(\{F^m:m\ge0\}\) is an Ore set. The first
basis is the PBW theorem with \(f_r\) last. In the localization, the finite
commutation formulas move any negative power of \(F\) to the right of every
ordered monomial. Linear independence follows by passing to the associated
graded: the associated graded of the localization is the localization of the
polynomial ring at the symbol \(\operatorname{gr}F\), so the exponent set
extends from \(\Z_{\ge0}\) to \(\Z\). Placing \(d\) first in the PBW order
gives the last assertions.
\end{proof}

\section{The boundary localization theorem}\label{sec:boundary}

Set
\[
F=f_b,\qquad E=e_{a+1},\qquad H=h_{a+b+1}.
\]
Define \(\psi:S_{a+1,b-1}\to\C\) by
\begin{align*}
\psi(K)=\varphi(K),\qquad
\psi(h_0)=\varphi(h_0)+2,\qquad
\psi(h_i)=\varphi(h_i)\ (i>0),
\end{align*}
and zero on the root vectors of \(S_{a+1,b-1}\). Thus \(\psi\) differs from
\(\varphi\) only in the \(h_0\)-eigenvalue.

\begin{lemma}[\(\psi\) is a character]\label{lem:char}
The displayed assignment defines a Lie algebra character
\(\psi:S_{a+1,b-1}\to\C\).
\end{lemma}

\begin{proof}
Only brackets that can fail are \([e_i,f_j]=h_{i+j}+i\delta_{i+j,0}K\) with
\(i>a+1\), \(j>b-1\). Then \(i+j\ge a+b+2=L+1\), so the central term vanishes
and \(\psi(h_{i+j})=\lambda_{i+j}=0\). The remaining brackets
\([h_i,e_j]=2e_{i+j}\), \([h_i,f_j]=-2f_{i+j}\) land in the annihilated tails,
and \([h_i,h_j]=2i\delta_{i+j,0}K\) with \(i,j\ge0\) is either zero or has
coefficient zero.
\end{proof}

\begin{lemma}[Cyclic vector]\label{lem:cyc}
The vector \(\bar u=F^{-1}u+M^c_{a,b}(\varphi)\in\T_FM^c_{a,b}(\varphi)\)
satisfies
\[
x\bar u=\psi(x)\bar u\qquad(x\in S_{a+1,b-1}).
\]
\end{lemma}

\begin{proof}
For \(j>b-1\): if \(j=b\), then
\(f_b\bar u=u+M^c_{a,b}(\varphi)=0\); if \(j>b\), then \(f_ju=0\). From
\[
h_iF^{-1}=F^{-1}h_i+2F^{-2}f_{i+b}
\]
one gets \(h_0\bar u=(\lambda_0+2)\bar u\), since
\(2F^{-2}f_bu=2F^{-2}Fu=2\bar u\), and \(h_i\bar u=\lambda_i\bar u\) for
\(i>0\), since \(f_{i+b}u=0\). For \(j>a+1\),
\[
e_jF^{-1}=F^{-1}e_j-F^{-2}h_{j+b}-2F^{-3}f_{j+2b};
\]
here \(j+b>a+b+1=L\), so \(h_{j+b}u=0\), and \(j+2b>b\), so \(f_{j+2b}u=0\).
Hence \(e_j\bar u=0\), and \(K\bar u=\kappa\bar u\).
\end{proof}

\begin{theorem}[Boundary localization]\label{thm:boundary}
If
\[
\lambda_\partial=\varphi(h_{a+b+1})\ne0,
\]
then there is a natural \(\slh\)-module isomorphism
\[
\boxed{\ \T_{f_b}M^c_{a,b}(\varphi)\cong M^c_{a+1,b-1}(\psi),\ }
\]
sending the standard generator \(u'\) of the target to
\[
\bar u=f_b^{-1}u+M^c_{a,b}(\varphi).
\]
The statement is independent of the level \(\kappa=\varphi(K)\).
\end{theorem}

\begin{proof}
By Lemma~\ref{lem:cyc} and the universal property of induction, there is a
unique \(\slh\)-module map
\[
\Phi:M^c_{a+1,b-1}(\psi)\longrightarrow\T_FM^c_{a,b}(\varphi),
\qquad u'\longmapsto\bar u.
\]
Compare PBW bases directly. Let \(\X\) be the ordered monomials in
\[
e_i\ (i\le a),\qquad h_j\ (j<0),\qquad f_k\ (k\le b-1),
\]
and choose the PBW order of the source so that \(E=e_{a+1}\) is the last free
variable. Then
\[
\{XE^pu':X\in\X,\ p\ge0\}
\]
is a basis of \(M^c_{a+1,b-1}(\psi)\), and by Lemma~\ref{lem:right-pbw}
\[
\{XF^{-p-1}u+M^c_{a,b}(\varphi):X\in\X,\ p\ge0\}
\]
is a basis of the localization quotient. The commutation relations give
\[
[E,F]=H,\qquad (\ad F)^2(E)=-2f_{a+2b+1}.
\]
Now \(Eu=0\), \(Hu=\lambda_\partial u\), and
\(a+2b+1=b+(a+b+1)>b\), so \(f_{a+2b+1}u=0\). Thus for every \(s\in\Z\),
\[
EF^su=s\lambda_\partial F^{s-1}u,
\]
and in particular
\[
E^pF^{-1}u=(-1)^pp!\lambda_\partial^{\,p}F^{-p-1}u.
\]
Therefore
\[
\Phi(XE^pu')
=(-1)^pp!\lambda_\partial^{\,p}XF^{-p-1}u+M^c_{a,b}(\varphi),
\]
and \(\Phi\) sends a PBW basis bijectively onto nonzero scalar multiples of a
PBW basis. Hence \(\Phi\) is an isomorphism.
\end{proof}

\begin{remark}[Surjection plus simplicity shortcut]\label{rem:shortcut}
One may avoid the base comparison when only simplicity is needed. The chain
\(EF^{-m}u=-m\lambda_\partial F^{-m-1}u\) shows that every pure fraction
\(F^{-m}u+M\) lies in the image of \(\Phi\); a numerator-degree induction as in
\cite[Lemma 2.2]{Nonboundary} shows that \(\Phi\) is surjective. If
\(\kappa\ne-2\), then by Theorem~\ref{thm:fgxz} the source
\(M^c_{a+1,b-1}(\psi)\) is simple, because
\[
\psi(h_{(a+1)+(b-1)+1})=\psi(h_L)=\lambda_\partial\ne0,\qquad
\psi(K)=\kappa\ne-2.
\]
A nonzero surjection from a simple module is an isomorphism, and the quotient
is simple.
\end{remark}

\begin{corollary}[Simplicity of the boundary quotient]\label{cor:simple}
If
\[
\lambda_\partial\ne0\qquad\text{and}\qquad \kappa\ne-2,
\]
then \(\T_{f_b}M^c_{a,b}(\varphi)\) is simple. More precisely, it is
isomorphic to \(M^c_{a+1,b-1}(\psi)\), whose two irreducibility parameters are
\[
\psi(h_{a+b+1})=\varphi(h_{a+b+1})=\lambda_\partial,
\qquad
\psi(K)=\varphi(K)=\kappa.
\]
\end{corollary}

\begin{remark}[Why not the full localization]\label{rem:full}
The full localization \(\D_FM^c_{a,b}(\varphi)\) contains
\(M^c_{a,b}(\varphi)\) as a nonzero \(\slh\)-submodule, so it cannot be simple
as soon as the localization is strictly larger. The object carrying the
standard-module structure is the quotient \(\T_FM=\D_FM/M\). This remains true
for twisted localizations, which keep \(F\) invertible.
\end{remark}

\section{Bernstein--Sato-type computation}\label{sec:bsato}

The chain used in Theorem~\ref{thm:boundary} is a rank-one \(b\)-function
computation in the enveloping-algebra localization, not a \(D_X\)-module
Bernstein--Sato polynomial.

\begin{proposition}[Boundary \(b\)-function]\label{prop:bfun}
With \(F=f_b\), \(E=e_{a+1}\) as above,
\[
EF^su=s\varphi(h_{a+b+1})F^{s-1}u\qquad(s\in\Z),
\]
equivalently
\[
EF^{s+1}u=(s+1)\varphi(h_{a+b+1})F^su,
\]
so the monic \(b\)-polynomial is
\[
b(s)=s+1,\qquad\text{or, with parameters,}\qquad
b_\varphi(s)=(s+1)\varphi(h_{a+b+1}).
\]
The condition \(\lambda_\partial\ne0\) is exactly the absence of obstruction
in the negative-power direction, and the recursion
\[
E^pF^{-1}u=(-1)^pp!\lambda_\partial^{\,p}F^{-p-1}u
\]
is the cyclicity of the pure fractions.
\end{proposition}

\begin{proof}
From \([E,F]=H\) and \((\ad F)^2(E)=-2f_{a+2b+1}\) with
\(Eu=0\), \(Hu=\lambda_\partial u\), \(f_{a+2b+1}u=0\), the finite expansion
\[
EF^s=F^sE+sF^{s-1}H-2\tbinom{s}{2}F^{s-2}f_{a+2b+1}
\]
gives the claim on \(u\). See also \cite[Section 5]{FGM} for the analogous
role of such resonances for free-field modules.
\end{proof}

\section{Arbitrary modes: the trichotomy}\label{sec:trichotomy}

\begin{proposition}[Arbitrary-mode trichotomy]\label{prop:trichotomy}
Let \(M=M^c_{a,b}(\varphi)\) and consider \(F=f_r\).
\begin{enumerate}[label=\textup{(\roman*)},leftmargin=2.8em]
\item If \(r>b\), then \(F\) is locally nilpotent on \(M\), hence
\(\D_FM=0\): there is no nontrivial localization.
\item If \(r=b\) and \(\varphi(h_{a+b+1})\ne0\), then
\[
\T_{f_b}M^c_{a,b}(\varphi)\cong M^c_{a+1,b-1}(\psi)
\]
by Theorem~\ref{thm:boundary}.
\item If \(r<b\), then \(\D_FM\ne0\), but \(\T_FM\) is not isomorphic to any
standard module \(M^c_{a',b'}(\chi)\): the set of locally nilpotent negative
root modes is \(\{r\}\cup\{j:j>b\}\), which has a gap, while every standard
module has a full upper tail \(\{j:j>b'\}\).
\end{enumerate}
The same trichotomy holds for \(M^f_{a,b}(\varphi)\). The deep study of case
(iii) is carried out in \cite{Nonboundary}; only the boundary cases (i) and
(ii) are needed below.
\end{proposition}

\begin{proof}
For \(r>b\), \(Fu=0\) and \(F^m(xu)=(\ad F)^m(x)u\) vanishes for \(m\gg0\)
by local nilpotence of \(\ad F\), so \(F\) is locally nilpotent on \(M\) and
inverting it kills everything. For \(r<b\), the localized PBW basis of
Lemma~\ref{lem:right-pbw} shows that \(F\) is locally nilpotent on the
quotient, while \(f_j\) for \(j\le b\), \(j\ne r\) acts as an injective free
polynomial variable; for \(j>b\) the local nilpotence on \(M\) is inherited.
\end{proof}

\section{Adjoining the derivation}\label{sec:derivation}

On the vector space \(M^f_{a,b}(\varphi)\cong\C[d]\otimes M^c_{a,b}(\varphi)\)
the \(\slt\)-action is the semidirect one: for \(x_r=x\otimes t^r\) and
\(p(d)\in\C[d]\),
\[
d\bigl(p(d)\otimes v\bigr)=dp(d)\otimes v,\qquad
x_r\bigl(p(d)\otimes v\bigr)=p(d-r)\otimes x_rv.
\]
In particular, for \(F=f_b\),
\[
F^{-1}\bigl(p(d)\otimes v\bigr)=p(d+b)\otimes F^{-1}v.
\]

\begin{proposition}[Localization commutes with adjoining \(d\)]
\label{prop:loc-induction}
Let \(X\) be an \(F=f_b\)-torsion-free \(\slh\)-module. Then there are natural
\(\slt\)-module isomorphisms
\[
\D_F\bigl(\Ind_{\slh}^{\slt}X\bigr)\cong\Ind_{\slh}^{\slt}(\D_FX),
\qquad
\boxed{\ \T_F\bigl(\Ind_{\slh}^{\slt}X\bigr)\cong\Ind_{\slh}^{\slt}(\T_FX).\ }
\]
As vector spaces the second isomorphism reads
\[
\T_F\bigl(\C[d]\otimes X\bigr)\cong\C[d]\otimes\T_FX.
\]
\end{proposition}

\begin{proof}
Set \(A=U(\slh)\), \(B=U(\slt)\). Since \([d,F]=bF\), the derivation \(\ad d\)
extends to \(\D_FA\), and PBW gives the natural normal-form isomorphism
\(\D_FB\cong B\otimes_A\D_FA\). Hence
\(\D_F(B\otimes_AX)\cong B\otimes_A\D_FX\). Since \(B\) is free as a right
\(A\)-module, \(B\otimes_A-\) is exact, and applying it to
\(0\to X\to\D_FX\to\T_FX\to0\) gives the quotient isomorphism.
\end{proof}

\begin{theorem}[Boundary localization with the derivation]\label{thm:boundary-f}
If \(\varphi(h_{a+b+1})\ne0\), then
\[
\boxed{\ \T_{f_b}M^f_{a,b}(\varphi)\cong M^f_{a+1,b-1}(\psi).\ }
\]
In the PBW normal form \(M^f\cong\C[d]\otimes M^c\), the isomorphism is termwise
\[
d^kXE^pu'\longmapsto
(-1)^pp!\lambda_\partial^{\,p}\,d^kXF^{-p-1}u+M^f_{a,b}(\varphi).
\]
The iterated chains of Section~\ref{sec:iterated} hold verbatim for \(M^f\).
\end{theorem}

\begin{proof}
By Proposition~\ref{prop:loc-induction} and Theorem~\ref{thm:boundary},
\begin{align*}
\T_{f_b}M^f_{a,b}(\varphi)
&\cong\Ind_{\slh}^{\slt}\bigl(\T_{f_b}M^c_{a,b}(\varphi)\bigr)\\
&\cong\Ind_{\slh}^{\slt}M^c_{a+1,b-1}(\psi)
=M^f_{a+1,b-1}(\psi).
\end{align*}
The explicit formula is obtained by adjoining the powers \(d^k\) on the left
of the base comparison in Theorem~\ref{thm:boundary}.
\end{proof}

\section{Iterated boundary localizations}\label{sec:iterated}

For \(k\ge1\) define the character \(\varphi^{(k)}\) by
\[
\varphi^{(k)}(K)=\varphi(K),\qquad
\varphi^{(k)}(h_0)=\varphi(h_0)+2k,\qquad
\varphi^{(k)}(h_i)=\varphi(h_i)\ (i>0).
\]

\begin{theorem}[Iterated boundary chain]\label{thm:iterated}
Assume \(\varphi(h_{a+b+1})\ne0\). Then, for every \(k\ge1\),
\begin{equation}\label{eq:iterated}
\boxed{\ \T_{f_{b-k+1}}\cdots\T_{f_{b-1}}\T_{f_b}
M^c_{a,b}(\varphi)\cong M^c_{a+k,b-k}(\varphi^{(k)}).\ }
\end{equation}
Here the rightmost functor acts first. The same nondegeneracy condition
\(\varphi(h_{a+b+1})\ne0\) is used at every step, because \(a+b\), the
boundary character value at \(h_{a+b+1}\), and the level \(\kappa\) are
preserved along the chain. In particular, for \(M^c_{-1,N}(\varphi)\) and
\(n\le N\), taking \(k=N-n+1\),
\begin{equation}\label{eq:iterated-special}
\boxed{\ \T_{f_n}\T_{f_{n+1}}\cdots\T_{f_N}
M^c_{-1,N}(\varphi)\cong M^c_{N-n,n-1}(\varphi^{(N-n+1)}).\ }
\end{equation}
\end{theorem}

\begin{proof}
Apply Theorem~\ref{thm:boundary} successively: after the first step the next
mode \(f_{b-1}\) is the boundary mode of \(M^c_{a+1,b-1}(\psi)\), and the
character has moved from \(\varphi\) to \(\varphi^{(1)}\); after \(k\) steps
the indices and the \(h_0\)-shift are as displayed. For
\eqref{eq:iterated-special}, apply \eqref{eq:iterated} to
\((a,b)=(-1,N)\).
\end{proof}

This is the correct way to reach an interior mode \(f_n\) inside the standard
family: move the boundary step by step, never skip the intermediate
\(f\)-modes. A direct interior localization leaves the standard family; see
Proposition~\ref{prop:trichotomy}(iii) and \cite{Nonboundary}.

\section{Non-critical simple flow with the Casimir}\label{sec:omega}

Assume \(\kappa\ne-2\). The Casimir element \(\Omega\) of \(\slt\) commutes
with \(U(\slt)\), and for a fixed nonzero polynomial \(P\in\C[\Omega]\) and
\(\mu\in\C\) the paper \cite{FGXZ} forms the central quotient; to simplify
notation we take the representative \(P(\Omega)=1\) and write
\[
L_{a,b}(\varphi;\mu):=M^f_{a,b}(\varphi)/(\Omega-\mu)M^f_{a,b}(\varphi).
\]

\begin{theorem}[Non-critical simple flow]\label{thm:omega-flow}
Assume
\[
\varphi(h_{a+b+1})\ne0,\qquad \varphi(K)\ne-2.
\]
Then \(f_b\) is injective on \(L_{a,b}(\varphi;\mu)\), and
\[
\boxed{\ \T_{f_b}L_{a,b}(\varphi;\mu)\cong L_{a+1,b-1}(\psi;\mu).\ }
\]
Both sides are simple \(\slt\)-modules, and the Casimir parameter \(\mu\) does
not change. More generally,
\[
\T_{f_{b-k+1}}\cdots\T_{f_b}L_{a,b}(\varphi;\mu)
\cong L_{a+k,b-k}(\varphi^{(k)};\mu).
\]
\end{theorem}

\begin{proof}
By \cite[formula (4.1)]{FGXZ}, at level \(\kappa\ne-2\) the restriction of
\(L_{a,b}(\varphi;\mu)\) to \(\slh\) is isomorphic to \(M^c_{a,b}(\varphi)\);
hence the boundary vector \(f_b\) remains injective. Since \(\Omega\) commutes
with \(F=f_b\), exactness of Ore localization gives
\[
\T_F\bigl(M^f/(\Omega-\mu)M^f\bigr)
\cong
\bigl(\T_FM^f\bigr)/(\Omega-\mu)\bigl(\T_FM^f\bigr).
\]
Applying Theorem~\ref{thm:boundary-f}, the right side becomes
\[
M^f_{a+1,b-1}(\psi)/(\Omega-\mu)M^f_{a+1,b-1}(\psi)
=L_{a+1,b-1}(\psi;\mu).
\]
The irreducibility parameters of \cite[Theorem 4.1]{FGXZ} are unchanged under
the boundary move, so both sides are simple. Iterate.
\end{proof}

\begin{remark}[Why one cannot stop at \(M^f\)]\label{rem:omega}
At non-critical level, \((\Omega-\mu)M^f\) is a proper submodule of \(M^f\),
so the raw module \(M^f\) and the raw quotient
\(\T_{f_b}M^f\cong M^f_{a+1,b-1}(\psi)\) are generally not simple. The complete
simple flow is
\[
\boxed{\ \text{boundary localization quotient}\ +\ \Omega\text{-central reduction}.\ }
\]
The two operations commute, and \(\mu\) is preserved throughout.
\end{remark}

\section{Critical level: the \(\theta\)-parameters are preserved}\label{sec:critical}

Throughout this section \(\kappa=\varphi(K)=-2\) and
\(q:=a+b+1=L\). For the reduced case \((a,b)=(-1,q)\), the critical central
modes used in \cite{FGXZ} are the Segal--Sugawara modes \(T(q-i)\).

\begin{proposition}[Spectral flow on the Sugawara modes]\label{prop:spectral}
Let \(\sigma_k\) be the automorphism
\[
\sigma_k(e_n)=e_{n+k},\qquad
\sigma_k(f_n)=f_{n-k},\qquad
\sigma_k(h_n)=h_n+kK\delta_{n,0},\qquad
\sigma_k(K)=K.
\]
In the completed enveloping algebra,
\[
\boxed{\ \sigma_k(T(m))
=T(m)+\frac{k}{2}(K+2)h_m+\frac{k^2}{4}K(K+2)\delta_{m,0}.\ }
\]
Consequently, at the critical level \(K=-2\),
\[
\boxed{\ \sigma_k(T(m))=T(m)\ }
\]
as operators on any smooth module.
\end{proposition}

\begin{proof}
Write \(T(m)\) as the combination \(\tfrac14(2,2,1)\) of the three
normal-ordered sums of \cite[formula (2.3)]{FGXZ}. For \(k>0\), spectral flow
moves the normal-ordering cut of the first sum
\(\sum_i:e_{-i}f_{m+i}:\) from \(i>0\) to \(i>-k\), producing
\(kh_m+\tfrac{k(k-1)}{2}K\delta_{m,0}\); the second sum
\(\sum_i:f_{-i}e_{m+i}:\) produces
\(kh_m+\tfrac{k(k+1)}{2}K\delta_{m,0}\); the third sum
\(\sum_i:h_{-i}h_{m+i}:\) produces
\(2kKh_m+k^2K^2\delta_{m,0}\). Inserting the coefficients
\(\tfrac14(2,2,1)\) gives the displayed formula. Negative \(k\) follows from
\(\sigma_{k+\ell}=\sigma_k\sigma_\ell\), and setting \(K=-2\) gives the
invariance.
\end{proof}

For the critical character \(\varphi(K)=-2\), define
\[
Z_i^{a,b}:=\sigma_{-a-1}\bigl(T(q-i)\bigr),\qquad i\in\N,
\]
and let \(I^{a,b}_{\boldsymbol\theta}\) be the ideal of the critical center
generated by \(Z_i^{a,b}-\theta_i\) (\(i\in\N\)). Put
\[
\overline M^{\,c}_{a,b}(\varphi;\boldsymbol\theta)
:=M^c_{a,b}(\varphi)/I^{a,b}_{\boldsymbol\theta}M^c_{a,b}(\varphi),
\]
and
\[
\overline M^{\,f}_{a,b}(\varphi;\boldsymbol\theta)
:=\Ind_{\slh}^{\slt}\overline M^{\,c}_{a,b}(\varphi;\boldsymbol\theta).
\]
By Proposition~\ref{prop:spectral}, at level \(-2\)
\[
Z_i^{a,b}=T(q-i)=Z_i^{a+1,b-1},
\]
so the boundary move does not change \(\boldsymbol\theta=(\theta_i)\).

\begin{theorem}[Critical simple flow]\label{thm:critical-flow}
Assume
\[
\varphi(K)=-2,\qquad \varphi(h_q)\ne0.
\]
Then \(f_b\) is injective on \(\overline M^{\,c}_{a,b}(\varphi;\boldsymbol\theta)\), and
\[
\boxed{\ \T_{f_b}\overline M^{\,c}_{a,b}(\varphi;\boldsymbol\theta)
\cong \overline M^{\,c}_{a+1,b-1}(\psi;\boldsymbol\theta),\ }
\]
both sides being simple \(\slh\)-modules. The same holds with \(f\) in place
of \(c\) throughout, both sides being simple \(\slt\)-modules.
\end{theorem}

\begin{proof}
The critical irreducibility theorem of \cite{FGXZ} (reduce to \((-1,q)\) and
twist back) makes \(\overline M^{\,c}_{a,b}(\varphi;\boldsymbol\theta)\)
simple. The \(F\)-torsion vectors of any module form a submodule because
\(\ad F\) is locally nilpotent on \(U(\slh)\). On the other hand, the critical
PBW basis of the paper contains the nonzero vectors \(f_q^{\,m}\bar u\);
twisting back by \(\sigma_{a+1}\) shows that \(f_b\) is not locally nilpotent
on the whole quotient. Hence the \(F\)-torsion submodule is zero and \(f_b\)
is injective.

Since each \(T(m)\) commutes with \(\slh\) and with \(F=f_b\), exactness of
Ore localization gives
\[
\T_F\!\left(M^c_{a,b}/I^{a,b}_{\boldsymbol\theta}M^c_{a,b}\right)
\cong
\frac{\T_FM^c_{a,b}}{I^{a,b}_{\boldsymbol\theta}\bigl(\T_FM^c_{a,b}\bigr)}.
\]
Applying Theorem~\ref{thm:boundary} and
\(Z_i^{a,b}=Z_i^{a+1,b-1}=T(q-i)\), the right side is exactly
\(\overline M^{\,c}_{a+1,b-1}(\psi;\boldsymbol\theta)\), which is simple by the
same irreducibility theorem. The \(f\)-version follows from
Proposition~\ref{prop:loc-induction}, with simplicity on both sides from
\cite[Theorem 4.2]{FGXZ}.
\end{proof}

Iterating,
\[
\T_{f_{b-k+1}}\cdots\T_{f_b}
\overline M^{\,\bullet}_{a,b}(\varphi;\boldsymbol\theta)
\cong
\overline M^{\,\bullet}_{a+k,b-k}(\varphi^{(k)};\boldsymbol\theta),
\qquad \bullet\in\{c,f\},
\]
with \(\boldsymbol\theta\) unchanged along the entire chain.

\section{Verification status}\label{sec:verify}

In the verification round of 2026-09-02 the key formulas underlying
Sections~\ref{sec:pbw}, \ref{sec:boundary}, \ref{sec:derivation}, and
\ref{sec:critical} were checked symbolically with three independent engines
(Python/SymPy, Mathematica 14.3, GAP 4.15.1), 76 checks in total, all
passing. These cover: the localization formulas \(xF^s=\sum_j(-1)^j\binom{s}{j}F^{s-j}(\ad F)^j(x)\)
as a consistent Laurent action for all \(s,t\in\Z\) (both sides are
polynomials of degree at most three in \(s,t\), so the symbolic identity is
exact), the closed forms used in Lemma~\ref{lem:cyc} and
Theorem~\ref{thm:boundary}, the \(\sigma_k\) automorphism on all defining
relations, the character property of \(\psi\), the cyclic-vector relations,
the chain \(E^pF^{-1}u=(-1)^pp!\lambda_\partial^{\,p}F^{-p-1}u\) for a grid
of \((a,b)\), the derivation shifts \(F^{-1}p(d)=p(d+b)F^{-1}\), and the
composition law plus \(K=-2\) invariance of the formula in
Proposition~\ref{prop:spectral}. No Lean/Coq formalization is claimed, and
the paper-level theorems (the FGXZ criterion, the critical irreducibility
statement, and the completed-center computations) remain external inputs.

\begingroup
\raggedright
\begin{thebibliography}{9}
\bibitem{FGXZ}
V.~Futorny, X.~Guo, Y.~Xue, and K.~Zhao,
\emph{Smooth representations of affine Kac--Moody algebras},
arXiv:2404.03855v2 (2024; journal version 2025).
\href{https://arxiv.org/html/2404.03855v2}{arXiv full text}.

\bibitem{FGM}
V.~Futorny, D.~Grantcharov, and R.~A.~Martins,
\emph{Localization of free field realizations of affine Lie algebras},
arXiv:1404.7148v2 (2015).
\href{https://arxiv.org/html/1404.7148v2}{arXiv full text}.

\bibitem{Nonboundary}
Project note, \emph{Nonboundary localization quotients of smooth affine
\(\mathfrak{sl}_2\)-modules} (2026-09-11),
\texttt{nonboundary\_localization\_results\_en.tex}.

\bibitem{LivingV8}
Project living note, \emph{Smooth representations of affine Kac--Moody
algebras: PBW bases, irreducibility, and localization tasks},
version v8 (2026-09-02),
\texttt{smooth\_kac\_moody\_living\_note\_v8\_cn.tex}.
\end{thebibliography}
\endgroup
\end{document}
